\section*{2. 高斯混合分布与参数估计（2 分）}

给定二维高斯混合概率密度函数
\[
p(\mathbf{x})
=
\sum_{k=1}^{3}
\pi_k \mathcal{N}(\mathbf{x}; \boldsymbol{\mu}_k, \boldsymbol{\Sigma}_k),
\qquad
\mathbf{x}=(x_1,x_2)^T \in \mathbb{R}^2,
\]
其中混合系数
\[
\pi_1 = 0.3, \qquad \pi_2 = 0.4, \qquad \pi_3 = 0.3,
\]
均值向量
\[
\boldsymbol{\mu}_1 =
\begin{bmatrix}
0\\
0
\end{bmatrix},
\qquad
\boldsymbol{\mu}_2 =
\begin{bmatrix}
3\\
0
\end{bmatrix},
\qquad
\boldsymbol{\mu}_3 =
\begin{bmatrix}
2\\
2
\end{bmatrix},
\]
协方差矩阵
\[
\boldsymbol{\Sigma}_1 =
\begin{bmatrix}
0.5 & 0\\
0 & 0.8
\end{bmatrix},
\qquad
\boldsymbol{\Sigma}_2 =
\begin{bmatrix}
0.8 & 0.3\\
0.3 & 0.5
\end{bmatrix},
\qquad
\boldsymbol{\Sigma}_3 =
\begin{bmatrix}
0.6 & -0.2\\
-0.2 & 0.6
\end{bmatrix}.
\]

请从该分布中生成 $1000$ 个样本，完成以下两道题目。

\begin{enumerate}
    \item 用 EM 算法对生成的样本进行聚类。设聚类数为 $3$，实现 EM 算法，估计高斯混合模型参数，包括混合系数、均值向量和协方差矩阵。可视化聚类结果和每个高斯成分的等高线。比较估计参数与真实参数的差异。（1 分）

    \item 基于生成的样本，用三种非参数方法进行密度估计：直方图方法、高斯核密度估计和 $K$ 近邻密度估计。分别改变区间（bin）宽度、核函数的方差 $H^2$ 和近邻数 $K$，在二维平面上可视化估计出的密度函数。（1 分）
\end{enumerate}